Des entiers sont placés aux sommets d'un polygone à $2n$ sommets.

Une opération permet d'ajouter $1$, ou de retrancher $1$, aux nombres situés
à deux sommets voisins.
\begin{enumerate}
  \item Montrer que la somme alternée des nombres est invariante.
  \item Montrer que si tous les nombres peuvent être rendus égaux, alors la
  somme alternée initiale est nulle.
  \item Prouver réciproquement que, si la somme alternée initiale est nulle,
  alors il est possible de rendre tous les nombres égaux.
\end{enumerate}
\emph{Indication :} procéder successivement autour du polygone en corrigeant
les sommets les uns après les autres.
